\chapter{Thầy Cũng Có Lúc Quên Mình Từng Là Trò}
\markboth{Thầy Cũng Có Lúc Quên Mình Từng Là Trò}{Thầy Cũng Có Lúc Quên Mình Từng Là Trò}

\section{Thầy Quên Cảm Giác Bị Gọi Bất Ngờ}
\begin{truyen}{Thầy Quên Cảm Giác Bị Gọi Bất Ngờ}{Teachers Forget What Sudden Exposure Feels Like}
\chuhoa{H}{ọc sinh:} ``Có khi thầy gọi em lên bảng nhanh quá nên em hoảng, rồi em càng ngu hơn bình thường.''
\textit{(Student: ``Sometimes you call me to the board so suddenly that I panic, and then I become even more stupid than usual.'')}

\teacherpair{Em nói vậy nghe hỗn nhưng không hẳn sai. Có những lúc người dạy quen đứng ở phía kiểm tra nên quên mất cảm giác bị nêu tên giữa đám đông nó căng thế nào. Gọi học sinh lên bảng là bình thường. Nhưng làm nhục sự lúng túng của học sinh thì không phải cách dạy tử tế.}{``What you said sounds rude, but it is not entirely wrong. Sometimes teachers grow so used to being the ones who test that they forget how tense it feels to be named in front of a crowd. Calling a student to the board is normal. Humiliating that student's confusion is not good teaching.''}
\end{truyen}

\section{Khi Kỷ Luật Bị Dùng Như Cơn Cáu}
\begin{truyen}{Khi Kỷ Luật Bị Dùng Như Cơn Cáu}{When Discipline Becomes Personal Irritation}
\chuhoa{H}{ọc sinh:} ``Có lúc em thấy thầy phạt vì bực chứ không phải vì muốn tụi em khá hơn.''
\textit{(Student: ``Sometimes it feels like teachers punish because they are annoyed, not because they want us to improve.'')}

\teacherpair{Câu đó đụng tự ái người lớn, nhưng đáng nghe. Kỷ luật cần công bằng và có mục tiêu. Nếu người dạy trút luôn cơn cáu vào hình phạt, thì học sinh không học được bài học nào ngoài việc sợ quyền lực.}{``That line hurts adult pride, but it deserves to be heard. Discipline must be fair and purposeful. If a teacher pours irritation directly into punishment, students learn nothing except fear of power.''}
\end{truyen}

\section{Người Dạy Cũng Phải Tự Soi}
\begin{truyen}{Người Dạy Cũng Phải Tự Soi}{Teachers Must Examine Themselves Too}
\chuhoa{H}{ọc sinh:} ``Tụi em bị nhắc soi lại mình suốt. Còn thầy cô thì sao?''
\textit{(Student: ``We are constantly told to reflect. What about teachers?'')}

\teacherpair{Thì cũng phải soi. Người dạy mà không tự soi sẽ rất nhanh biến thói quen của mình thành chân lý. Còn khi đã tưởng mình luôn đúng, người ta dễ quên mất cảm giác non nớt, vụng về và sĩ diện của tuổi học trò.}{``Then teachers must reflect too. A teacher who never examines themselves quickly turns habits into truth. And once a person assumes they are always right, they easily forget the awkwardness, insecurity, and pride of student life.''}
\end{truyen}

\begin{baihoc}
Muốn dạy tốt không chỉ cần giỏi chuyên môn, mà còn cần giữ lại ký ức mình từng là học trò.
\textit{(To teach well requires not only expertise but also remembering what it felt like to be a student.)}
\end{baihoc}

\begin{ghinhoanh}
\textbf{Short English to remember:}

Authority needs memory.

Good teaching includes self-reflection.
\end{ghinhoanh}

\begin{tuhocnhanh}
1. Nếu là giáo viên, em có dễ quên cảm giác run của học sinh không? \textit{(If you were a teacher, would you easily forget how nervous students feel?)}

2. Khi có quyền nhắc người khác, em có tự nhắc mình không? \textit{(When you have the right to correct others, do you also correct yourself?)}
\end{tuhocnhanh}

\ngancach
